\documentclass{article}
\usepackage{graphicx}
\usepackage{amsmath}
\usepackage{amsfonts}
\usepackage{bookmark}
\usepackage{booktabs}
\usepackage{exercises}
\graphicspath{ {./images/} }
\usepackage{indentfirst}
\usepackage{amsthm}
\usepackage{yhmath}
\usepackage{placeins}
\newcommand*{\oldneg}{\mathord{\sim}}
\newcommand\showdiv[1]{\overline{\smash{)}#1}}
\newcommand\Mydiv[1]{%
\smash{\raise.3ex\hbox{$\big)$}}\mkern-8mu
        \overline{\enspace\strut#1}}
\setlength{\tabcolsep}{5pt}
\theoremstyle{definition}
\newtheorem{definition}{Definition}

\setlength{\parindent}{.5in}
\setlength{\parskip}{12pt}
\renewcommand{\arraystretch}{2}
\newcommand{\?}{\stackrel{?}{=}}

\title{Standardized Test (pp. 298-299)}
\begin{document}
\maketitle

\section*{1}
C

\section*{2}
If the interest is compounded daily, then it is compounded 365 times in a year. The model to represent this situation is:
\[ A=1000(1+\dfrac{0.03}{365})^{365t} \] meaning, after two years:
\[ A=1000(1+\dfrac{0.03}{365})^{365\times 2}\approx 1061.83 \]
The model to model the situation in which interest is compounded monthly (12 times a year) is:
\[ A=1000(1+\dfrac{0.03}{12})^{12t} \] meaning, after two years:
\[ A=1000(1+\drac{0.03}{12)^{12\times 2}\approx 1061.76 \]
meaning the difference is about $\$.07$.
\newline A

\section*{3}
B

\section*{4}
D

\section*{5}
C

\section*{6}
From the given information, we know that a model for the population is:
\[ A=290(1.009)^t \] where $t$ is measured in years. Thus, because if we let $t=7$ (the year is 2010):
\[ A=290(1.009)^7\approx 309 \] million people.
\newline B

\section*{7}
A

\section*{8}
D

\section*{9}
\[ y=e^{2x}-3 \] \[ x=e^{2y}-3 \] \[ x+3=e^{2y} \] \[ 2y=\ln\left(x+3\right) \] \[ y=\dfrac{\ln\left(x+3\right)}{2} \]
C

\section*{10}
\[ \log\left(x+3\right)+\log x=1 \] \[ \log\left(x+3\right)=1-\logx \] \[ x+3=10^{1-\logx} \]
\[ x+3=10\times 10^{-\log x} \] \[ x+3=10\times(10^{\logx})^{-1} \] \[ x+3=10\times x^{-1} \]
\[ x+3=\frac{10}{x} \] \[ x^2+3x=10 \] \[ (x+5)(x-2)=0 \] \[ x=-5\vee x=2 \]
Check for extraneous solutions:
\[ \log\left(-5+3\right)+\log -5\?1 \] which is impossible as we cannot take logarithms of negative numbers.
\[ \log\left(2+3\right)+\log 2=1 \] meaning the only solution is $x=2$.\newline B

\section*{11}
B

\section*{12}
A model matching the given equation is:
\[ A=500e^{0.035t} \]
meaning, after one year:
\[ A=500e^{0.035}\approx \$517.81 \]

\section*{13}
We are given the model:
\[ y=7.7e^{0.14x} \] and wish to find the value of $x$ when $y=50$. Thus:
\[ 50=7.7e^{0.14x} \] \[ \frac{50}{7.7}=e^{0.14x} \] \[ \ln\frac{50}{7.7}=0.14x \]
\[ 0.14x\approx 1.871 \] \[ x\approx 13.364 \]

\section*{14}
\[ 9^{2x+1}=3^{5x-1} \] \[ (3^2)^{2x+1}=3^{5x-1} \] \[ 3^{4x+2}=3^{5x-1} \]
\[ 4x+2=5x-1 \] \[ x=3 \]

\section*{15}
\[ \log_2 72-2\log_2 3=\log_2 72-\log_2 3^2 \] \[ =\log_2 72-\log_2 9 \] \[ =\log_2\frac{72}{9} \]
\[ =\log_2 8=3 \]

\section*{16}
We know that:
\[ 8=ab^0 \] and \[ 2=ab^2 \] meaning \[ a=8 \] and thus:
\[ 2=8b^2 \] \[ b^2=\frac{1}{4} \] \[ b=\pm\frac{1}{2} \] Thus, the model exponential function corresponding to these points is:
\[ y=8(\pm\frac{1}{2})^x \] meaning $f(5)$ is:
\[ 8(\pm\frac{1}{2})^5=\frac{1}{4}\vee -\frac{1}{4} \]

\section*{17}
\[ \log_3\left(5x+3\right)=5 \] \[ 5x+3=3^5 \]
\[ 5x+3=243 \] \[ x=48 \]
Check for extraneous solutions:
\[ \log_3\left(5(48)+3\right)\? 5 \] \[ \log_3 253=5 \]

\section*{18}
We know that the model for the tempereature $T$ after $t$ minutes is:
\[ T=(T_0-T_r)e^{-rt} \] where $T_r$ is the surrounding temperature, $T_0$ is the initial temperature, and $r$ is the rate of cooling.
Thus, the model for out situation is:
\[ 110=(232-68)e^{-20r} \] \[ \frac{110}{164}=e^{-20r} \] \[ \ln\frac{55}{82}=-20r \]
\[ -20r\approx -0.399 \] \[ r\approx 0.01995 \]
One way that you could increase the cooling rate would be by decreasing the room temperature. This would make $T_0-T_r$ a lesser amount, therefore 
decreasing the total temperature.

\section*{19}
We know that $\$37$ million is the initial amount, and the percent decrease is $0.3$, meaning the model is:
\[ A=37(0.7)^t \] where $t$ is measured in weeks. First, the domain must be $x\ge 0$, as we cannot have negative values for $t$. Going in the positive 
direction, unless we know what the ticket price of the movie is, we can't make a reasonable assumption about the domain of $x$ in the positive $x$-axis, 
other than it must end at some finite number.

\section*{22}
(a): At 10,000 feet, the air density is about:
\[ 0.0761e^{-0.0004\times 10000}\approx 0.001 \]


(b): We wish to find the height at which the air pressure is 12 pounds per square inch:
\[ 12=14.7e^{-0.0004h} \] \[ \frac{12}{14.7}=e^{-0.0004h} \] \[ \ln\frac{12}{14.7}=-0.0004h \]
\[ -0.0004h\approx -0.203 \] \[ h\approx 507.5 \]


(c): At 60 degrees F, the air pressure is is greater than the air density by a factor of:
\[ \frac{14.7}{0.0761}\approx 193.167 \]


(d): We know that there are about 0.3048 meters in a foot, meaning this hypothesis would mean that there would be a $1%$ decrease every 
$80\times 0.3048=24.384$ feet. So, suppose we have a given air pressure of:
\[ 14.7e^{-0.0004h} \] at height $h$ feet, meaning, after the next $24.384$ feet up, the air pressure would be:
\[ 14.7e^{-0.0004(h+24.384)}=14.7e^{-0.0004h-0.0097392}=14.7e^{-0.0004h}\times e^{-0.0097392} \]
\[ \approx 0.99\times 14.7e^{-0.0004h} \]
Thus, the air density decreases by about $1%$, confirming the hypothesis.

\end{document}

